This section focuses on some of the more important concepts in ProcControlAPI
and gives a high level overview before the detailed API is presented in Section.

\subsection{Processes and Threads}
There are two central classes to ProcControlAPI, Process and Thread.
Each class respectively represents a single
target process or thread running on the system. By performing operations on
the Process and Thread objects, a ProcControlAPI user is able to control the
target process and its threads.

Each Process is guaranteed to have at least one Thread associated with it. A
multi-threaded process may have a Process object with more than one Thread.
Each process has an address space associated with it, which can be written or
read through the Process object. Each thread has a set of registers
associated with it, which can be access through the Thread object. 

At any one time a Thread will be in either a stopped state or a running
state. A thread in a
stopped state has had its execution paused by
ProcControlAPI- the OS will not schedule the thread to run.
A thread in a running state is allowed to execute as normal. A thread in a
running state may block for other reasons, e.g. blocking on IO calls, but this
does not affect ProcControlAPI\'s view of the thread state.
A thread is only in the stopped state if ProcControlAPI has explicitly
stopped it. 

A Process object is not considered to have a stopped or running
state- only its Thread objects are stopped or running. A
stop operation on a Process triggers a stop operation on each of its Threads,
and similarly a continue operation on a Process triggers continue operations on
each Thread.

\subsection{Callbacks}
In addition to controlling a target process through the Process and Thread
objects, a ProcControlAPI user can also receive notification of events that
happen in that process. Examples of these events would be a new thread being
created, a breakpoint being executed, or a process exiting.

The ProcControlAPI user receives notice of events through a callback system.
The user can register callback function that will be called by ProcControlAPI
whenever a particular type of event occurs. Details about the event are
passed to the callback function via an Event object.

\subsubsection{Events}
Each event can be broken up into
an EventType object and an Event
object. The EventType describes a type of event that can happen, and Event
describes a specific instance of an event happening. Each Event will have one
and only one EventType. 

Each EventType has two primary fields: its time and its code. The code field
of describes what type of event occurred, e.g. EventType::Exit represents a
target process exiting. The time field of an EventType represents whether the
EventType is happening before or after will have code and will have a value of
EventType::Pre, EventType::Post, or EventType::None. 

For example, an EventType with time and code of EventType::Pre and
EventType::Exit will occur just before a target process exits, and a code of
EventType::Exec with a time of EventType::Post will occur after an exec system
call occurs. In this document we will abbreviate EventTypes such as these as
pre-exit and post-exec. Some EventTypes do not have a time associated with
them, for example EventType::Breakpoint does not have an associated time and
thus has a time value of EventType::none. 

An Event represents an instance of an EventType occurring. In addition to an
EventType, each Event also has pointer to the Process and Thread that it
occurred on. Certain events may also have event specific information
associated with them, which is represented in a sub-class of Event. Each
EventType is associated with a specific sub-class of Event. 

For example, EventType::Library is used to signify a shared library being loaded
into the target process. When an EventType::Library occurs ProcControlAPI
will deliver an object of type EventLibrary, which is a subclass of Event, to
any registered callback functions. In addition to the information inherited
from Event, the EventLibrary will contain extra information about the library
that was loaded into the target process.

Table 1 shows the Event subclass that is used for each EventType. Not all
EventTypes are available on every platform- a checkmark
under the specific OS column means that the EventType is available on that OS.

\begin{table}[!ht]
  \centering
  \begin{tabular}{|l|l|l|l|l|}
  \hline
    EventType & Event Subclass & Linux & FreeBSD & Windows \\ \hline
    Stop & EventStop & \checkmark & ~ & ~ \\ \hline
    Breakpoint & EventBreakpoint & \checkmark & \checkmark & \checkmark \\ \hline
    Signal & EventSignal & \checkmark & \checkmark & \checkmark \\ \hline
    UserThreadCreate & EventNewUserThread & \checkmark & \checkmark & ~ \\ \hline
    LWPCreate & EventNewLWP & \checkmark & ~ & \checkmark \\ \hline
    Pre-UserThreadDestroy & EventUserThreadDestroy & \checkmark & \checkmark & ~ \\ \hline
    Post-UserThreadDestroy & EventUserThreadDestroy & \checkmark & \checkmark & ~ \\ \hline
    Pre-LWPDestroy & EventLWPDestroy & \checkmark & ~ & \checkmark \\ \hline
    Post-LWPDestroy & EventLWPDestroy & \checkmark & ~ & ~ \\ \hline
    Pre-Fork & EventFork & ~ & ~ & ~ \\ \hline
    Post-Fork & EventFork & \checkmark & ~ & ~ \\ \hline
    Pre-Exec & EventExec & ~ & ~ & ~ \\ \hline
    Post-Exec & EventExec & \checkmark & \checkmark & ~ \\ \hline
    RPC & EventRPC & \checkmark & \checkmark & \checkmark \\ \hline
    SingleStep & EventSingleStep & \checkmark & \checkmark & \checkmark \\ \hline
    Breakpoint & EventBreakpoint & \checkmark & \checkmark & \checkmark \\ \hline
    Library & EventLibrary & \checkmark & \checkmark & \checkmark \\ \hline
    Pre-Exit & EventExit & \checkmark & ~ & ~ \\ \hline
    Post-Exit & EventExit & \checkmark & \checkmark & \checkmark \\ \hline
    Crash & EventCrash & \checkmark & \checkmark & \checkmark \\ \hline
    ForceTerminate & EventForceTerminate & \checkmark & \checkmark & \checkmark \\ \hline
  \end{tabular}
\end{table}


\subsubsection{Callback Functions}
Events are delivered via a
callback function. A ProcControlAPI user can register callback functions for
an EventType using the Process::registerEventCallback function. All callback
functions must be declared using the signature:

\begin{lstlisting}[language=c++]
Process::cb_ret_t callback_func_name(Event::ptr)
\end{lstlisting}

In order to prevent a class of race conditions, ProcControlAPI does not allow a
callback function to perform any operation that would require another callback
to be recursively delivered. At most one callback function can be running at
a time.

To enforce this, the event that is passed to a callback function contains only
const pointers to the triggering Process and Thread objects. Any member
function that could trigger callbacks is not marked const, thus triggering a
compilation error if they are called on an object passed to a callback. If
the ProcControlAPI user uses const\_cast or global variables to get around the
const restriction it will result in a runtime error. API functions that
cannot be used from a callback are mentioned in the API entries.

Operations such as Process::stopProc, Process::continueProc, Thread::stopThread, and
Thread::continueThread are not safe to
call from a callback function, but it would still be useful to perform these
operations. ProcControlAPI allows the user to use the return value from a
callback function to specify whether process or thread that triggered the event
should be stopped or continued. More details on this can be found in the
\code{Process::cb\_ret\_t} section of the API reference.

\subsubsection{Callback Delivery}
When ProcControlAPI needs to
deliver a callback it must first gain control of a user visible thread in the
controller process. This thread will be used to invoke the callback
function. ProcControlAPI does not use its
internal threads for delivering callbacks, as this would expose the
ProcControlAPI user to race conditions. 

Unfortunately, the user thread is not always accessible to ProcControlAPI when
it needs to invoke a callback function. For example, the user visible thread
may be performing network IO or waiting for input from a GUI when an event
occurs. 

ProcControlAPI uses a notification system built around the EventNotify class to
alert the ProcControlAPI user that a callback is ready to be delivered. Once
the user is notified then they can call the Process::handleEvents function,
under which ProcControlAPI will invoke any pending callback functions. 

The EventNotify class has two mechanisms for
notifying the ProcControlAPI user that a callback is pending: writing to a file
descriptor and a light-weight callback function. The
EventNotify::getFD function returns a file
descriptor that will have a byte written to it when a callback is ready. This
file descriptor can be added to a select or poll to block a thread that handles
ProcControlAPI events. Alternatively, the ProcControlAPI user can register a
light-weight callback that is invoked when a callback is ready. This
light-weight callback provides no information about the Event and may occur on
another thread or from a signal handler- the ProcControlAPI
user is encouraged to keep this callback minimal. 

It is important for a user to respond promptly to a callback notification. A
target process may remain blocked while a notification is pending. If a
target process is generating many events that need callbacks, a long delay in
notification could have a significant performance impact.

Once the ProcControlAPI user knows that a callback is ready to be delivered they
can call Process::handleEvents, which will
invoke all callback functions. Alternatively, if the ProcControlAPI user does
not need to handle events outside of ProcControlAPI, they can continue to block
in Process::handleEvents without going through the notification system.

\subsection{Inferior Remote Procedure Call}
An iRPC is a mechanism for executing code in a target process. Despite the
name, an iRPC does not necessarily have to involve a procedure
call- any piece of code can be executed.

A ProcControlAPI user can invoke an iRPC by providing ProcControlAPI with a
buffer of machine code and specifying a Process or Thread on which to run the
machine code. ProcControlAPI will insert the machine code into the address
space, save the register set, run the machine code, and then remove the machine
code after execution completes. When the iRPC completes (but before the
registers and memory are cleaned) ProcControlAPI will deliver an
EventIRPC to any registered callback function. The
ProcControlAPI user may use this callback to collect any results from the
registers or memory used by the iRPC.

Note that ProcControlAPI will preserve the registers of the thread running the
iRPC, and it will preserve the memory used by the machine code. Other memory
or system state changed by the iRPC may remain visible to the target process
after the iRPC completes.

The machine code for each iRPC must contain at least one trap instruction (e.g.,
a 0xCC instruction on x86 family or a 0x7D821008 instruction on the PPC
family). ProcControlAPI will stop executing the iRPC upon invocation of the
trap. Note that the trap instruction must fall within the original machine
code for the iRPC. If the iRPC calls or jumps to another piece of code that
executes a trap instruction then ProcControlAPI will not treat it as the end of
the iRPC.

Before an iRPC can be run it must be posted to a process or thread using the
Process::postIRPC or
Thread::postIRPC API functions. The
Process::postIRPC function will select a thread to post the iRPC to.
Multiple iRPCs can be posted to the same thread, but only one iRPC will run
at a time- subsequent iRPCs will be queued and run after
the preceding iRPC completes. If multiple iRPCs are posted to different
threads in a multi-threaded process, then they may run in parallel.

An iRPC can be posted to a stopped or running thread. If posted to a stopped
thread, then the iRPC will run when the thread is continued. If posted to a
running thread, then the iRPC will run immediately or, if posted from a
callback function, when the callback function completes.

An iRPC may be blocking or non-blocking. If a blocking iRPC is posted to any
Process, then calls to Process::handleEvents will block until the iRPC is completed.
